% ===================================================================
%  GAMBAR No. 2 — Awak Manungsa. --- Wektu Ngaso. Dolanan.
%
%  Traduction, faite en 2026, en javanais d'aujourd'hui, sur l'ido de
%  texto/io. Regles de la colonne : voir 10-gambar-01.tex. Une seule
%  numerotation.
%
%  LE CORPS HUMAIN EST LE TABLEAU OU LE JAVANAIS NE DOIT RIEN A
%  PERSONNE, ET LE SEUL DU LIVRET QUI SOIT DANS CE CAS. Pas un mot de
%  neerlandais, pas un de portugais, pas un d'arabe : sirah la tete,
%  rai le visage, bathuk le front, mripat l'oeil, irung le nez,
%  cangkem la bouche, lambe les levres, untu les dents, ilat la
%  langue, gulu le cou, githok la nuque, dhadha la poitrine, geger le
%  dos, pundhak l'epaule, weteng le ventre, lengen le bras, sikut le
%  coude, ugel-ugel le poignet, driji les doigts, getih le sang,
%  dhengkul le genou, kempol le mollet, polok la cheville, balung
%  l'os, tlapakan la plante. Une langue emprunte les choses qu'on lui
%  apporte ; le corps, personne ne le lui a apporte.
%
%  ET LA DIFFERENCE AVEC L'INDONESIEN Y EST TOTALE : sirah contre
%  kepala, rai contre wajah, mripat contre mata, irung contre hidung,
%  untu contre gigi, ilat contre lidah, gulu contre leher, geger
%  contre punggung, pundhak contre bahu, weteng contre perut, getih
%  contre darah, dhengkul contre lutut, balung contre tulang. Deux
%  langues voisines a se confondre partout ailleurs, et pas un mot
%  commun sur trente parties du corps. Le lexique le plus ancien est
%  celui qui se ressemble le moins.
%
%  LES CINQ SENS SONT CINQ NOMS VERBAUX D'UN SEUL PREFIXE.
%  pandeleng, pangrungu, pangambu, pangecap, pangrasa : voir,
%  entendre, sentir, gouter, toucher, chacun pris a son verbe par le
%  meme « pa- ». La ou l'ido compose cinq fois « -senso » et le
%  francais decline cinq noms sans parente, le javanais fait une
%  serie, et elle se lit comme telle. C'est le seul endroit du volume
%  ou la morphologie de la langue rend la lecon mieux que la lecon.
%
%  MAIS cangkem EST UN MOT DE NIVEAU, ET IL FALLAIT LE CHOISIR. La
%  bouche se dit cangkem en ngoko et tutuk en krama ; le nez irung en
%  ngoko et grana en krama ; manger mangan en ngoko et dhahar en
%  krama inggil. La colonne tient le ngoko (voir l'en-tete du tableau
%  1), donc cangkem, irung, mangan — et non parce que ce serait plus
%  simple : parce que c'est un enfant qui raconte, et qu'un enfant
%  javanais qui parlerait de son propre corps en krama se moquerait
%  de lui-meme.
%
%  L'ECHELLE (2) TRANCHE LE FAUX AMI ANNONCE AU TABLEAU 1. On y
%  disait que « tangga » est le voisin en javanais et l'echelle en
%  indonesien. Le voici en situation : l'echelle du portique se dit
%  andha, et « tangga » ne pourrait pas la nommer sans dire autre
%  chose. Le controle ne releve toujours pas le mot — il est juste
%  dans les deux emplois javanais —, mais le tableau montre pourquoi
%  il fallait s'en mefier.
%
%  LE VELO (37) SE DIT pit, ET C'EST LE SECOND COUPLE APRES
%  potlot/pensil. Le javanais l'appelle pit, du neerlandais fiets ;
%  l'indonesien l'appelle sepeda, du neerlandais velocipede. Meme
%  langue preteuse, meme decennie, meme machine, et deux mots qui
%  n'ont pas une lettre en commun. Ce n'est pas que l'une ait emprunte
%  et l'autre non : c'est qu'elles n'ont pas emprunte le meme mot au
%  meme Hollandais. Et neker (21), la bille, vient du meme cote —
%  knikker — quand l'indonesien dit kelereng.
%
%  LA NOTE (*) DU TABLEAU AUTORISE LE NOM NATIONAL, ET LA COLONNE LE
%  REFUSE DEUX FOIS. Guignon ecrit que les jeux d'enfants « havas
%  ofte nomi pure nacionala » et qu'il a pris le nom national « kande
%  lu ne esas tro nejusta ». Le javanais en offrait un tout pret pour
%  le jeu de barres (38) : gobak sodor, qui se joue encore dans toutes
%  les cours de Java. Il n'a pas ete pris — « dolanan palang » —
%  parce que la planche grave une ronde francaise et non un gobak
%  sodor, et que lui donner ce nom-la aurait modernise LA CHOSE. Le
%  bilboquet (19), a l'inverse, n'a aucun nom javanais et garde le
%  sien : bilboket. On emprunte le nom quand la chose est venue avec ;
%  on ne pose pas un nom d'ici sur une chose de la-bas. La source est
%  ici plus permissive que la colonne, et c'est note parce que c'est
%  un desaccord, non un oubli.
%
%  UNE FAUTE DE FORME, RELEVEE PAR LA MACHINE : « \newline » N'EXISTE
%  PAS DANS CE RELEVE. La note (*) porte deux alineas, et l'ido les
%  separe avec ses \nl de fin de ligne — que la traduction n'a pas.
%  On avait ecrit \newline pour les separer ; kolonoj.py l'a releve
%  comme macro inconnue des la premiere passe. La colonne indonesienne
%  avait deja resolu la meme difficulte, et sans macro : les deux
%  alineas se suivent dans les memes accolades, separes par un retour
%  a la ligne du fichier. Le controle des macros connues ne sert pas
%  qu'a attraper les fautes de frappe — il empeche d'inventer une
%  mise en page que le volume n'a pas.
%
%  ET LE BLOC t02-25-1 A ETE REPRIS DE LA MEME FACON. Le fac-simile
%  ido le coupe en deux par un feuillet, la note s'intercalant ; on
%  avait recopie cette coupure avec un « p suite ». Une traduction
%  n'a ni feuillet ni folio : le bloc est ecrit d'un seul tenant et
%  la note vient apres, comme en indonesien. L'ordre des renvois est
%  le meme dans les deux dispositions — 24, 25, 28, 29 —, mais la
%  seconde est la seule qui obeisse a la regle de la colonne.
%
%  ET LES JEUX QUI ONT UN NOM JAVANAIS LE GARDENT TOUS : gangsingan
%  la toupie (18), egrang les echasses (24), dhelikan la cache-cache
%  (25), layangan le cerf-volant (33), bal-balan le football (36),
%  kodhok la grenouille du jeu de rainette (20). Aucun n'a ete
%  traduit ni glose : ce sont les memes jeux, et ils ont ici leur
%  nom.
% ===================================================================

%%K t02-apar-1 apar
\VUsaut{4.0mm}
\VUtitre{15.0pt}{60}{GAMBAR No. 2}
\VUsaut{2.0mm}
\VUtitre{11.4pt}{0}{Awak Manungsa. --- Wektu Ngaso.}
\VUtitre{11.4pt}{0}{Dolanan.}

%%K t02-tit-0 sub
\VUtitre{13.2pt}{0}{Awak Manungsa.}
\VUsaut{2.4mm}
\VUcentre{10.2pt}{0}{\textit{(Piwulang prasaja bab ngelmu alam.)}}
\VUsaut{3.0mm}

%%K t02-tit-1 sub pos=t02-01-1
\VUpk{10.2pt}{40}{Sirah.}
\VUsaut{3.0mm}

%%K t02-01-1 p
1. --- Perangan baku awak manungsa yaiku : \VUgras{sirah,} \VUgras{badan}
lan \VUgras{anggota awak.}

%%K t02-02-1 p
2. --- Sirah kaperang dadi loro : \VUgras{bathok sirah,} kang isi
\VUgras{utek} lan katutupan \VUgras{rambut,} sarta \VUgras{rai.}

%%K t02-03-1 p
3. --- Ing gambar kita ora weruh bathok sirah utawa utek ; nanging
kita weruh cetha banget perangan-perangan wigati ing rai.

%%K t02-04-1 p
4. --- Iki kang kapisan \VUgras{bathuk,} kang nuduhake pikiran landhep,
kang tansah lumaku. \VUgras{Pilingane} amba banget. \VUgras{Mripate} urip
lan mlolo. \VUgras{Alis} kandel lan \VUgras{idep} dawa ngayomi mripat iku.

%%K t02-05-1 p
5. --- Iki uga \VUgras{irung} lan \VUgras{bolongan irunge.} Irung migunani
kanggo ngambu lan ambegan. Ing saben sisihe irung ana \VUgras{pipi,}
luwih adoh maneh \VUgras{kuping.} Perangan ngisore rai yaiku
\VUgras{cangkem} lan \VUgras{janggut.}

%%K t02-06-1 p
6. --- Kita ora weruh cetha banget, amarga \VUgras{brengos} lan
\VUgras{jambang} ngalingi. Nanging kita ngerti yen cangkem duwe
\VUgras{lambe} loro, \VUgras{rahang ndhuwur} lan \VUgras{rahang ngisor}
karo \VUgras{untune,} sarta \VUgras{ilat.} Nganggo cangkem kita guneman
lan mangan. Nganggo ilat kita ngrasakake pangan.

%%K t02-07-1 p
7. --- Mripat, kuping, irung lan cangkem, padha karo driji, yaiku
piranti limang pancadriya : \VUgras{pandeleng,} \VUgras{pangrungu,}
\VUgras{pangambu,} \VUgras{pangecap,} \VUgras{pangrasa.}

%%K t02-08-1 p
8. --- Dadi sirah iku salah sijine perangan awak manungsa kang paling
narik kawigaten.

%%K t02-tit-2 sub
\VUsaut{4.4mm}
\VUpk{10.2pt}{40}{Badan.}
\VUsaut{3.0mm}

%%K t02-09-1 p
9. --- \VUgras{Gulu} nyambungake sirah karo badan. Gulu kalebu
\VUgras{gorokan} lan \VUgras{githok.}

%%K t02-10-1 p
10. --- Badan uga isi akeh piranti kang wigati. Ing \VUgras{dhadha} ana
\VUgras{paru-paru,} kang kanggo ambegan, lan \VUgras{jantung} kang dadi
puseran urip. Yen jantung mandheg thing-thing, titah kang urip iku
mati.

%%K t02-11-1 p
11. --- \VUgras{Iga} nyangga dhadha.

%%K t02-12-1 p
12. --- Perangan liyane badan yaiku \VUgras{iringan,} \VUgras{geger,}
\VUgras{pundhak} lan \VUgras{weteng.} Kita turu miring ing iringan utawa
nglumah ing geger, lan kita manggul momotan ing pundhak.

%%K t02-13-1 p
13. --- Ing weteng ana \VUgras{lambung} lan \VUgras{usus.} Loro-lorone
migunani kanggo ngluluhake pangan.

%%K t02-tit-3 sub
\VUsaut{4.4mm}
\VUpk{10.2pt}{40}{Anggota Awak.}
\VUsaut{3.0mm}

%%K t02-14-1 p
14. --- Kita duwe \VUgras{anggota awak} papat : \VUgras{lengen} loro lan
\VUgras{sikil} loro. Nganggo lengen kita njupuk, nyekel, ngangkat lan
nglumpukake barang ; nganggo sikil kita ngadeg, mlaku lan mlayu.

%%K t02-15-1 p
15. --- \VUgras{Pundhak} nyambungake lengen karo awak. \VUgras{Otot}
kang kuwat banget, yaiku bisep, nggerakake lengen kang disambung
\VUgras{sikut} menyang \VUgras{lengen ngisor.} \VUgras{Ugel-ugel} dadi
sendhine \VUgras{tangan,} kang pungkasane \VUgras{driji} lan
\VUgras{kuku.} Ing tangan kita weruh \VUgras{urat getih} — lan uga
nadi — kang diliwati \VUgras{getih.}

%%K t02-16-1 p
16. --- Perangane sikil yaiku : \VUgras{pupu,} \VUgras{dhengkul} lan
\VUgras{mburine dhengkul,} \VUgras{kempol} kang \VUgras{daginge}
katutupan \VUgras{kulit,} \VUgras{polok} lan \VUgras{tlapakan sikil.}
\VUgras{Balung} sikil gedhe banget lan kuwat banget. Ing pucuking sikil
ana \VUgras{driji sikil.} Perangan liyane yaiku \VUgras{tapak sikil}
lan \VUgras{tungkak.}

%%K t02-17-1 p
17. --- Mangkono gegambaran awak manungsa kang meh jangkep.

%%K t02-17-2 p
\textit{Cathetan.} --- \textit{Kanthi ukara « aku lara ing... (sirah lan
sapanunggalane) » guru bisa nganggo maneh sakabehe tembung iki saka
sudut pandang becik utawa alane} \VUgras{kaanan awak.}

%%K t02-tit-4 sub
\VUsaut{5.0mm}
\VUpk{10.2pt}{40}{Senam.}
\VUsaut{2.4mm}
\VUcentre{10.2pt}{0}{\textit{(Critane salah sijine murid.)}}
\VUsaut{3.0mm}

%%K t02-18-1 p
18. --- Supaya lentur, trampil lan waras, kita kudu nglatih sikil lan
lengen. Mula kita dolanan lan nglakoni \VUgras{senam.}

%%K t02-19-1 p
19. --- Sajrone seminggu, latihan loro iku dianakake ing latar sekolah
menengah (delengen gambar No. 1). Nanging ing dina Kemis lan Minggu
kita menyang palemahan suket kang jembar, panggonan kita duwe papan
kanggo mlayu lan seneng-seneng.

%%K t02-20-1 p
20. --- Guru-guru lan wong tuwa kita kadhang ngeterake kita mrana ;
nanging \VUgras{guru senam}\textsuperscript{(12)} kang nuntun latihan iku.

%%K t02-21-1 p
21. --- Ing palemahan suket, ing sisih kiwane lawang mlebu, ana
\VUgras{piranti senam}\textsuperscript{(1)} ; kita menek nganggo
\VUgras{andha}\textsuperscript{(2)}, kita malik awak ing \VUgras{gelang-gelang
senam}\textsuperscript{(3)} lan ing \VUgras{palang ayun}\textsuperscript{(4)},
kita mendhat awak nganggo tangan tekan pucuking \VUgras{andha
tali}\textsuperscript{(5)} utawa \VUgras{tali simpulan}\textsuperscript{(6)}.

%%K t02-22-1 p
22. --- Sauntara iku, kanca-kanca kita mlumpat ngliwati \VUgras{jaran
kayu}\textsuperscript{(7)} utawa \VUgras{palang sejajar}\textsuperscript{(11)}.
Liyane \VUgras{tinju}\textsuperscript{(8)} utawa \VUgras{anggar}\textsuperscript{(9)}
nganggo kedhok lan \VUgras{floret.} Pungkasane, liyane maneh
\VUgras{ngangkat halter}\textsuperscript{(10)}.

%%K t02-tit-5 sub
\VUsaut{4.6mm}
\VUpk{10.2pt}{40}{Dolanan.}
\VUsaut{3.0mm}

%%K t02-23-1 p
23. --- Ing sakubenge bocah-bocah kang senam, bocah lanang lan bocah
wadon padha dolanan warna-warna \VUgras{dolanan.} Bocah wadon
seneng-seneng nganggo \VUgras{gelang gedhene}\textsuperscript{(14)} ; padha
mlumpat nganggo \VUgras{tali}\textsuperscript{(13)} utawa ngajak adhi lan
sedulur lanange main \VUgras{kroket}\textsuperscript{(15)}. Padha bungah
banget bisa nyabetake \VUgras{bal} mungsuhe adoh nganggo \VUgras{palu
kayune,} pas nalika mungsuhe ngira bakal gampang liwat lengkungan!

%%K t02-24-1 p
24. --- Bocah lanang luwih seneng dolanan kang mbutuhake otot. Padha
seneng main \VUgras{kegel}\textsuperscript{(16)} kang dirubuhake kanthi
trampil nganggo \VUgras{bal}\textsuperscript{(17)}. Padha ora nampik uga
dolanan \VUgras{gangsingan}\textsuperscript{(18)} lan
\VUgras{bilboket}\textsuperscript{(19)} utawa malah \VUgras{dolanan
kodhok}\textsuperscript{(20)}. Nanging dolanan kang paling disenengi yaiku
\VUgras{neker}\textsuperscript{(21)} lan \VUgras{bal tembok}\textsuperscript{(22)},
amarga temboke \VUgras{omah kaca}\textsuperscript{(26)} pancen mathuk banget
kanggo mbalekake bal kang entheng iku.

%%K t02-25-1 p
25. --- Ioannes cilik, kang mlaku nganggo \VUgras{egrang}\textsuperscript{(24)},
trampil banget lan bisa njaga imbang kanthi becik. Ing cedhake,
sawetara kanca dolanan \VUgras{dhelikan}\textsuperscript{(25)} ing omah
kaca iku lan ing mburine \VUgras{pager.} Liyane luwih seneng
\VUgras{dolanan bedhekan}\textsuperscript{(28)} utawa
\VUgras{badminton}\textsuperscript{(29)}, kang koke mabur saka
\VUgras{raket} siji menyang liyane.

%%K t02-noto-1 noto
\VUnotes{143.2mm}{%
(*) Dolanan bocah asring duwe jeneng kang murni saka bangsane dhewe.
Kita ngusahakake menehi jeneng ing basa Ido kanthi becik-becike, ana
kalane kanthi ngupaya saka tumindake dhewe kang dadi ciri dolanan iku,
ana kalane kanthi milih jeneng bangsa ing negara iki utawa negara kana,
yen jeneng mau ora banget-banget kliru. Tuladha : \VUgras{barelo-ludo}
(Jerman lan Prancis) yaiku \VUgras{rano-ludo} \textsuperscript{(20)}
(Inggris), kang ing kono kang dolanan ngupaya nguncalake cakram
bunderane liwat cangkeme kodhok wesi kang menga ngadeg ing kothak
(barelo) kang bolong.
\VUgras{Shultro-salto}\textsuperscript{(30)} beda karo \VUgras{dorso-salto}
mung ing bab iki : kanggo mlumpat ngliwati sirah, kang mlumpat nyanggong
tangane ing pundhake kancane, dene ing « \VUgras{dorso-salto} » tangane
disanggongake ing gegere thok. --- Utawa uga sawetara kang melu padha
mbungkuk dene liyane bakal mlumpat ngliwati gegere kanthi nyanggong
tangan ing pundhak ; kang kapisan kudu tahan kena tumbukan tanpa
ambruk, kang kapindho kudu mlumpat tanpa ilang imbange (Delengen
gambare dhewe).%
}

%%K t02-26-1 p
26. --- Ing tengahe palemahan suket ana wit loro. Ing sangisore salah
sijine, bocah lanang pitu utawa wolu nganakake \VUgras{lumpat
pundhak}\textsuperscript{(30)} (*). Ora ing kabeh negara dolanan iku
dingerteni ; nanging kabeh wong ngerti apa iku \VUgras{lumpat geger,}
kang ora didolanake bocah-bocah nalika iku.

%%K t02-tit-6 sub
\VUsaut{5.0mm}
\VUpk{10.2pt}{40}{Dolanan ing Jaba.}
\VUsaut{3.4mm}

%%K t02-27-1 p
27. \VUgras{Dolanan picekan}\textsuperscript{(31)} uga nyenengake banget ;
bocah wadon lan bocah lanang gentenan masang \VUgras{tutup mata} ing
mripate banjur nyekel kang kurang trampil, kang nganti kena dicekel.

%%K t02-28-1 p
28. --- Ing tengahe palemahan suket, iki dolanan kang Prancis banget
nanging rada kasar : yaiku \VUgras{dolanan bruwang}\textsuperscript{(32)},
kang luwih rame tinimbang dolanan \VUgras{layangan}\textsuperscript{(33)}
kang tentrem banget iku, utawa malah tinimbang dolanan Inggris
\VUgras{tenis}\textsuperscript{(34)}.

%%K t02-29-1 p
29. --- Olahraga kang keri iku dadi bungahe murid-murid gedhe.
Guru-guru kita dhewe seneng nglakoni. Olahraga iku njaluk katrampilan
gedhe lan kacukatan, lan aku seneng banget. Dolanan
\VUgras{ayunan}\textsuperscript{(35)} dakpasrahake marang bocah wadon.

%%K t02-30-1 p
30. --- Aku ora seneng dolanan Inggris liyane kang bisa uga dadi
mbebayani.

%%K t02-30-1b p
Kang dakkarepake yaiku \VUgras{bal-balan}\textsuperscript{(36)} kang
mbungahake kangmasku. Aku luwih seneng \VUgras{dolanan
palang}\textsuperscript{(38)} lawas kita, kang padha becike kanggo
kasarasan lan pancen kurang kasar.

%%K t02-tit-7 sub
\VUsaut{4.6mm}
\VUpk{10.2pt}{40}{Numpak pit lan dolanan balon.}
\VUsaut{3.0mm}

%%K t02-31-1 p
31. --- Aku uga seneng ngubengi palemahan suket nganggo
\VUgras{pit}\textsuperscript{(37)}.

%%K t02-32-1 p
32. --- Taun ngarep aku bakal sinau \VUgras{nunggang
jaran}\textsuperscript{(39)}. Endah temen jaran kang mlaku utawa nyongklang
kanthi luwes, lan kang mlumpat ngliwati alangan!

%%K t02-33-1 p
33. --- Ing mangsa panas, kita sinau \VUgras{nglangi}\textsuperscript{(40)}.
Kita seneng nyilem lan adus ing banyu kali kang bening lan seger.
Kadhang ing pungkasan kita nguncalake \VUgras{balon}\textsuperscript{(41)}
kertas kang munggah kanthi agung menyang awang-awang, sanalika balon
iku kebak hawa panas.

%%K t02-34-1 p
34. --- Isih akeh dolanan liyane, nanging aku ora bakal rampung
ngetung kabeh, lan saka ringkesan iki wong ngerti yen ing dina Kemis,
wong ora bosen ing palemahan suket kita kang endah.

%%K t02-34-2 p
N.-B. --- \textit{Guru bakal gampang njangkepi bab iki ; kanthi
nggambarake luwih rinci saben dolanan kang paling wigati.}
